
\documentclass[../../../physics-standard_answers_all.tex]{subfiles}

\begin{document}

\setlength{\abovedisplayskip}{2mm}
\setlength{\belowdisplayskip}{1mm}

\begin{enumerate}

    \item まず，$z_1 = \uwave{0}$ であり，キルヒホッフ則より，
        %
        \begin{align*}
            \frac{x_1}{C_0} = V_0
            \qquad \text{∴} \; x_1 = \uwave{C_0 V_0} \;.
        \end{align*}
        %
        すると，
        %
        \begin{align*}
            E_1 = \frac{x_1}{C_0 \cdot 2d} = \uwave{\frac{V_0}{2d}} \;.
        \end{align*}

    \item XとZからなる部分についての電荷保存則より，
        %
        \begin{align*}
            x_2 + z_2 = x_1 + z_1
            \qquad \text{∴} \; x_2 + z_2 = C_0 V_0 \;.
        \end{align*}
        %
        キルヒホッフ則より，
        %
        \begin{align*}
            \frac{x_2}{C_0} = \frac{z_2}{2C_0} 
            \qquad \text{∴} \; x_2 = \uwave{\frac{1}{3} C_0 V_0} \;,
            \quad z_2 = \uwave{\frac{2}{3} C_0 V_0} \;.
        \end{align*}
        %
        すると，
        %
        \begin{align*}
            E_2 = \frac{x_2}{C_0 \cdot 2d} = \uwave{\frac{V_0}{6d}} \;,
            \quad \phi_2 = \frac{x_2}{C_0} = \uwave{\frac{1}{3} V_0} \;.
        \end{align*}

    \item 電荷保存則より，Zの電荷は変わらない．つまり，
        %
        \begin{align*}
            z_3 = z_2 = \uwave{\frac{2}{3} C_0 V_0} \;.
        \end{align*}
        %
        キルヒホッフ則より，
        %
        \begin{align*}
            \frac{x_3}{C_0} = V_0 
            \qquad \text{∴} \; x_3 = \uwave{C_0 V_0} \;.
        \end{align*}

    \item キルヒホッフ則より，
        %
        \begin{align*}
            \frac{x_4}{C_0} = \frac{z_4}{2C_0} = V_0
            \qquad \text{∴} \; x_4 = \uwave{C_0 V_0} \;,
            \quad z_4 = \uwave{2C_0 V_0} \;.
        \end{align*}

    \item XとZからなる部分についての電荷保存則より，
        %
        \begin{align*}
            x_5 + z_5 = x_4 + z_4
            \qquad \text{∴} \; x_5 + z_5 = 3C_0 V_0 \;.
        \end{align*}
        %
        キルヒホッフ則より，
        %
        \begin{align*}
            &\frac{x_5}{C_0} = \frac{z_5}{\varepsilon_r \cdot 2C_0} \\
            & \text{∴} \; x_5 = \uwave{\frac{3}{2\varepsilon_r + 1} C_0 V_0} \;,
            \quad z_5 = \uwave{\frac{6\varepsilon_r}{2\varepsilon_r + 1} C_0 V_0} \;.
        \end{align*}
        
\end{enumerate}

\end{document}
